
\subsection{Intermittent radial transport: finite variation and exact axis-jet torque}
\label{sec:xli-intermittent}

The dissipation obstruction of Edition~XL excludes sustained growth of
raw radial angular-momentum flux at the parabolic shrinking radius.
The remaining intermittent mechanism must therefore exploit radial
differentiation or sharply localized space--time flux profiles.
We derive a quantitative bounded-variation budget, retain the angular
Fourier leakage \emph{after} differentiation, and compute the leading
axis coefficient of the Reynolds torque from the Cartesian velocity jet.
All results are necessary conditions for a single actual unforced
Navier--Stokes flow, not a construction of an unforced blowup solution.

\begin{theorem}[Space--time bounded variation of the nonlinear radial flux]
\label{thm:xli-BV}
Let $u$ be a smooth finite-energy exactly unforced Navier--Stokes
solution in $\mathbb R^3$ with viscosity $\nu>0$ and
$E_0=\|u(0)\|_2^2$. Choose $\eta\in C_c^\infty(\mathbb R)$
supported in an interval of length $L$, and define
\[
F_\rho(t)=\rho^2\int_{\mathbb R}\int_0^{2\pi}
\eta(z)u_r(\rho,\theta,z,t)u_\theta(\rho,\theta,z,t)\,
d\theta dz.
\]
For any time interval $I$ before the singular time, put
$D_I=\int_I\|\nabla u(t)\|_2^2dt$, and take any valid
Sobolev constant $S_3$ in
$\|v\|_6\le S_3\|\nabla v\|_2$ on $\mathbb R^3$.
Then
\begin{equation}\label{eq:xli-BV}
\boxed{\begin{aligned}
\int_I\sup_{0<\rho\le R}|F_\rho(t)|dt
&\le\int_I\int_0^R|\partial_\rho F_\rho(t)|d\rho dt\\
&\le \|\eta\|_\infty D_I
\bigl[S_3^2(\pi L)^{2/3}R^{4/3}
 +S_3(\pi L)^{1/3}R^{5/3}\bigr]\\
&\le\frac{\|\eta\|_\infty E_0}{2\nu}
\bigl[S_3^2(\pi L)^{2/3}R^{4/3}
 +S_3(\pi L)^{1/3}R^{5/3}\bigr].
\end{aligned}}
\end{equation}
For the terminal interval $I_\tau=[1-2\tau,1-\tau]$
and $R_\tau=2\sqrt{X\tau}$ the left sides are
$o(\tau^{2/3})$, provided the smooth unforced
flow exists on $[0,1)$ with finite initial energy.
This is an integrated radial-variation estimate and
does not exclude a spike at one radius or instant.
\end{theorem}
\begin{proof}
Differentiate the actual flux:
\[
\partial_\rho F_\rho=
2\rho\int\eta u_ru_\theta+
\rho^2\int\eta[(\partial_\rho u_r)u_\theta+
u_r(\partial_\rho u_\theta)].
\]
Both integrals are in $(\theta,z)$.
Since $2|u_ru_\theta|\le |u|^2$ and
$|(\partial_\rho u_r)u_\theta+
u_r\partial_\rho u_\theta|\le
|u||\partial_\rho u|$,
\[
\int_I\int_0^R|\partial_\rho F_\rho|d\rho dt
\le \|\eta\|_\infty
(A_{I,R}+R\sqrt{A_{I,R}B_{I,R}}),
\]
where $A_{I,R}=\int_I\int_{\{r<R,z\in\mathrm{supp}\eta\}}
|u|^2dxdt$ and $B_{I,R}$ is the same integral
with $|\nabla u|^2$. The cylinder volume is
$\pi R^2L$. H\"older, the whole-space Sobolev
inequality and $B_{I,R}\le D_I$ give
$A_{I,R}\le S_3^2(\pi R^2L)^{2/3}D_I$.
Smooth axis compatibility gives $F_0=0$;
radial total variation thus dominates its supremum.
Finally the exact unforced dissipation inequality
$2\nu\int_0^1\|\nabla u\|_2^2dt\le E_0$
proves the uniform budget. Since
$D_{I_\tau}\to0$ by absolute continuity of
this finite integral, the terminal estimate
improves to little-$o$.
\end{proof}

\begin{corollary}[Intermittent derivative spike cost]
\label{cor:xli-spikes}
If $W_\tau\subset I_\tau$ and
$J_\tau\subset(0,R_\tau)$ have measures
$|W_\tau|$ and $\ell_\tau$ and
$|\partial_\rho F_\rho(t)|\ge A_\tau$ almost
everywhere on $W_\tau\times J_\tau$, then
\begin{equation}\label{eq:xli-spike}
\boxed{
A_\tau |W_\tau|\ell_\tau\le
\|\eta\|_\infty D_{I_\tau}
\bigl[S_3^2(\pi L)^{2/3}R_\tau^{4/3}
+S_3(\pi L)^{1/3}R_\tau^{5/3}\bigr].
}
\end{equation}
For a time window of length comparable to $\tau$,
radial width comparable to $\sqrt\tau$, and
amplitude comparable to $\tau^{-\alpha}$,
a fixed positive lower bound is impossible for
$\alpha\ge5/6$. Smaller exponents are not
excluded by this estimate.
\end{corollary}
\begin{proof}
Integrate the hypothesized absolute derivative
over the stated rectangular set and apply
Theorem~\ref{thm:xli-BV}. Compare
$\tau^{3/2-\alpha}$ with $o(\tau^{2/3})$.
\end{proof}

\begin{proposition}[Angular leakage after taking the flux derivative]
\label{prop:xli-leakage}
For $r>0$ define the angular fluctuation covariances
$J=\overline{u'_r\Gamma'}$ and
$K=\overline{u'_z\Gamma'}$, where $\Gamma=ru_\theta$.
Truncate their Fourier series to
$1\le |m|\le N$ and let
$\mathcal T_N=-r^{-1}\partial_r(rJ_N)-\partial_zK_N$.
Then
\begin{equation}\label{eq:xli-leakage}
\boxed{\begin{aligned}
|\mathcal T-\mathcal T_N|
\le\frac1{2\pi(N+1)^2}\bigl[
&r^{-1}U_\theta G_\theta+
U_{\theta r}G_\theta+U_\theta G_{\theta r}\\
&+Z_{\theta z}G_\theta+Z_\theta G_{\theta z}
\bigr],
\end{aligned}}
\end{equation}
where $U_\theta=\|\partial_\theta u_r\|_{L^2_\theta}$,
$Z_\theta=\|\partial_\theta u_z\|_{L^2_\theta}$,
$G_\theta=\|\partial_\theta\Gamma\|_{L^2_\theta}$,
and subscripts $r,z$ indicate the corresponding
mixed angular-radial or angular-axial derivative norms,
all at fixed $(r,z,t)$.
\end{proposition}
\begin{proof}
The covariance is the exact sum over $m\ne0$
of products of angular Fourier coefficients.
Differentiate each discarded product. Apply
Cauchy--Schwarz and Parseval with
$|m|\ge N+1$. The $r^{-1}J$ contribution
remains explicit near the symmetry axis.
\end{proof}

\begin{proposition}[Leading axis-jet coefficient of nonlinear torque]
\label{prop:xli-axis-jet}
Let $U=u_x,V=u_y,W=u_z$ denote Cartesian
velocity components, and evaluate the following
quantities and spatial derivatives at $(0,0,z,t)$:
$a=U(0,0,z,t)$, $b=V(0,0,z,t)$, and
\[
J_3=\tfrac18[a(-2U_{xy}+V_{xx}-V_{yy})
+b(U_{xx}-U_{yy}+2V_{xy})],\qquad
K_2=\tfrac12(bW_x-aW_y).
\]
For a sufficiently smooth field, as $r\downarrow0$,
\begin{equation}\label{eq:xli-axis-jet}
\boxed{
J(r,z,t)=r^3J_3+O(r^5),\quad
K(r,z,t)=r^2K_2+O(r^4),\quad
\mathcal T(r,z,t)=r^2[-4J_3-\partial_zK_2]+O(r^4).
}
\end{equation}
All purely first-transverse-derivative products
cancel from the leading radial covariance $J_3$.
For the actual exact unforced Gaussian datum
from Edition~XXXVII, the coefficient equals
$-4\varepsilon^2e^{-2z^2}$, in agreement
with its full initial torque.
\end{proposition}
\begin{proof}
Taylor expand $U,V$ transversely through second
order and $W$ through first order.
Substitute $x=r\cos\theta,y=r\sin\theta$
in $u_r=U\cos\theta+V\sin\theta$ and
$\Gamma=xV-yU$, integrate all trigonometric
products, and subtract the products of angular
means. The first-derivative covariance terms
cancel exactly. The angular mean enforces
the indicated parity of radial powers;
controlled higher Cartesian jets give the
remainders. Differentiate the covariance
according to
$\mathcal T=-r^{-1}\partial_r(rJ)-\partial_zK$.
Independent exact symbolic Gaussian substitution
agrees with the result.
\end{proof}

\paragraph{Proof boundary.}
At the source similarity radius $r_\tau\asymp\sqrt\tau$,
bounded Cartesian jets suppress the torque by $r_\tau^2$.
Source-compatible positive growth would require
amplifying jets or the controlled Taylor remainder
within \emph{one} exactly unforced evolution.
Pressure has no direct mean angular torque, but
its nonlocal effect on velocity cannot be omitted.
The original published Section~10 Navier--Stokes
background uses external forcing. Neither
the required positive terminal torque lower bound,
a singularity-preserving correction, nor a full
Lean PDE formalization has been proved.
